\documentclass[journal]{IEEEtran}
\usepackage[utf8]{inputenc}
\usepackage[spanish,es-tabla]{babel}
\usepackage{graphicx,amsmath,booktabs,array,url}
\usepackage[hidelinks]{hyperref}

% Marcador visible para lo que falta por completar/verificar
\newcommand{\TODO}[1]{\textbf{[COMPLETAR: #1]}}

\begin{document}

\title{PrediCasa: Predicción del precio de viviendas mediante regresión supervisada}
\author{\IEEEauthorblockN{Roller Andrés Hernández López, Juan Miguel Cadena Zúñiga y Juan Felipe Martínez Patiño}\\
\IEEEauthorblockA{Departamento de Ingeniería de Sistemas, Universidad de Antioquia, Colombia\\
Modelos y Simulación de Sistemas II}}
\maketitle

\begin{abstract}
Este informe presenta la primera entrega del proyecto PrediCasa, cuyo objetivo es estimar el precio de venta de viviendas a partir de sus características físicas, de calidad y de ubicación, usando técnicas de aprendizaje de máquina. Se describe el problema y su contexto, la composición del conjunto de datos \textit{Ames Housing} (competencia \textit{House Prices: Advanced Regression Techniques} de Kaggle), la configuración de aprendizaje seleccionada (aprendizaje supervisado, regresión) y una revisión de trabajos previos sobre la misma base de datos.
\end{abstract}

\begin{IEEEkeywords}
predicción de precios, vivienda, regresión, aprendizaje supervisado, Ames Housing.
\end{IEEEkeywords}

\section{Descripción del problema}

\subsection{Contexto y utilidad de una solución basada en ML}
La fijación del precio de una vivienda suele hacerse a partir de la impresión subjetiva de compradores, vendedores y agentes inmobiliarios, o de comparaciones superficiales como los metros cuadrados o el número de habitaciones. Sin embargo, el valor de un inmueble depende de muchos otros factores: la calidad de construcción y de los acabados, la ubicación (barrio), la presencia de garaje, sótano o chimenea, el año de construcción y si ha sido remodelado, entre otros. Ignorar estas variables puede llevar a sobrevalorar o subvalorar una propiedad, generando pérdidas para vendedores, dificultando la compra a los compradores y distorsionando el mercado.

Un modelo de aprendizaje de máquina puede aprender de ventas históricas cómo se combinan decenas de variables (incluyendo relaciones no lineales e interacciones) para producir una estimación objetiva, reproducible y cuantificable del precio, y además permite identificar qué características aportan más al valor. PrediCasa busca ofrecer esa estimación como apoyo a la toma de decisiones, no como reemplazo del criterio experto.

\subsection{Composición de la base de datos}
Se usa el conjunto \textit{Ames Housing}, compilado por De Cock \cite{decock2011} y publicado en Kaggle \cite{kaggle}, que contiene ventas de viviendas residenciales en Ames (Iowa, EE.\,UU.) entre 2006 y 2010. La versión de la competencia se divide en:
\begin{itemize}
\item \textbf{Entrenamiento:} 1460 muestras y 81 columnas (Id, 79 variables explicativas y la variable objetivo \texttt{SalePrice}).
\item \textbf{Prueba:} 1459 muestras y 80 columnas, sin \texttt{SalePrice} (etiquetas ocultas para la evaluación en Kaggle).
\end{itemize}
Las 79 variables describen la vivienda y el lote. Ejemplos: \texttt{OverallQual} (calidad general de materiales y acabados, escala 1--10), \texttt{GrLivArea} (área habitable sobre el suelo, en pies$^2$), \texttt{Neighborhood} (barrio), \texttt{YearBuilt} y \texttt{YearRemodAdd} (año de construcción y de remodelación), \texttt{GarageCars} (capacidad del garaje), \texttt{TotalBsmtSF} (área del sótano). La variable objetivo, \texttt{SalePrice}, es el precio de venta en dólares (variable continua, con asimetría positiva). En el conjunto de entrenamiento, el precio medio es USD\,180\,921 (desviación estándar 79\,443), la mediana USD\,163\,000, el mínimo USD\,34\,900 y el máximo USD\,755\,000. La asimetría es 1.88 y la curtosis 6.54: la distribución tiene una cola larga hacia las viviendas costosas, por lo que se trabajará con $\log(1+y)$. El conjunto tiene 38 columnas numéricas (incluyendo \texttt{Id} y \texttt{SalePrice}) y 43 categóricas, sin filas duplicadas.

\textbf{Datos faltantes.} Varias columnas presentan valores nulos; según la documentación del conjunto, en la mayoría de categóricas el valor \texttt{NA} no significa dato perdido sino \textit{ausencia de la característica} (por ejemplo, sin piscina, sin callejón, sin sótano, sin garaje). Las más afectadas son \texttt{PoolQC}, \texttt{MiscFeature}, \texttt{Alley}, \texttt{Fence} y \texttt{FireplaceQu}, seguidas de \texttt{LotFrontage} y las variables de garaje y sótano. En total, 19 columnas presentan nulos (Tabla~\ref{tab:nulos}). Los datos realmente perdidos son \texttt{LotFrontage} (259 casos), \texttt{MasVnrType}/\texttt{MasVnrArea} (8 casos) y \texttt{Electrical} (1 caso).

\begin{table}[!t]
\caption{Variables con datos faltantes (1460 viviendas).}
\label{tab:nulos}
\centering\small
\begin{tabular}{lrr}
\toprule
Variable(s) & Nulos & \% \\
\midrule
\texttt{PoolQC} & 1453 & 99.52 \\
\texttt{MiscFeature} & 1406 & 96.30 \\
\texttt{Alley} & 1369 & 93.77 \\
\texttt{Fence} & 1179 & 80.75 \\
\texttt{FireplaceQu} & 690 & 47.26 \\
\texttt{LotFrontage} & 259 & 17.74 \\
\texttt{Garage*} (5 variables) & 81 & 5.55 \\
\texttt{BsmtExposure}, \texttt{BsmtFinType2} & 38 & 2.60 \\
\texttt{BsmtQual}, \texttt{BsmtCond}, \texttt{BsmtFinType1} & 37 & 2.53 \\
\texttt{MasVnrType}, \texttt{MasVnrArea} & 8 & 0.55 \\
\texttt{Electrical} & 1 & 0.07 \\
\bottomrule
\end{tabular}
\end{table}

\textbf{Estrategia de imputación.} (i) En categóricas donde \texttt{NA} indica ausencia (y en \texttt{MasVnrType}), se reemplaza por la categoría ``None''; (ii) en numéricas asociadas (\texttt{GarageYrBlt}, \texttt{MasVnrArea}) se imputa 0; (iii) \texttt{LotFrontage} se imputa con la mediana por barrio (\texttt{Neighborhood}); (iv) los pocos casos restantes (p.\,ej. \texttt{Electrical}) se imputan con la moda. Todas las estadísticas de imputación se calculan solo con el conjunto de entrenamiento para evitar fuga de información. Tras la imputación no quedan valores nulos.

\textbf{Codificación de variables.} Las variables se agrupan en: numéricas continuas y discretas (se mantienen, con posible transformación logarítmica de las asimétricas y de \texttt{SalePrice}); ordinales de calidad/condición (\texttt{ExterQual}, \texttt{KitchenQual}, \texttt{BsmtQual}, etc.), codificadas con enteros que respetan el orden (Po$<$Fa$<$TA$<$Gd$<$Ex); y nominales (\texttt{Neighborhood}, \texttt{MSZoning}, \texttt{HouseStyle}, etc.), codificadas con \textit{one-hot}. Variables numéricas que en realidad son categóricas (p.\,ej. \texttt{MSSubClass}, \texttt{MoSold}) se tratan como nominales. Las variables numéricas se estandarizan para los modelos sensibles a la escala. En total se codificaron 10 variables ordinales y 36 nominales; tras el \textit{one-hot} el conjunto queda con 255 predictores (matriz de 1460$\times$256 incluyendo \texttt{SalePrice}).

\subsection{Hallazgos del análisis exploratorio}
La Fig.~\ref{fig:corr} muestra las variables numéricas más correlacionadas con el precio: la calidad general (\texttt{OverallQual}, $r=0.79$) y el área habitable (\texttt{GrLivArea}, $r=0.71$) encabezan la lista, seguidas de la capacidad y el área del garaje ($r=0.64$ y $0.62$) y del área del sótano ($r=0.61$). Esto confirma que el precio depende de mucho más que los metros cuadrados o el número de habitaciones. Entre los predictores hay pares con correlación superior a 0.8 (\texttt{GarageCars}--\texttt{GarageArea}, 0.88; \texttt{YearBuilt}--\texttt{GarageYrBlt}, 0.83; \texttt{GrLivArea}--\texttt{TotRmsAbvGrd}, 0.83; \texttt{TotalBsmtSF}--\texttt{1stFlrSF}, 0.82), lo que indica multicolinealidad que se tratará en el modelado (selección o combinación de variables y regularización).

\begin{figure}[!t]
\centering
\includegraphics[width=\columnwidth]{fig_corr.png}
\caption{Top 15 de variables numéricas más correlacionadas con \texttt{SalePrice}.}
\label{fig:corr}
\end{figure}

Las viviendas con garaje tienen una mediana de precio de USD\,167\,500 frente a USD\,100\,000 sin garaje, y las que tienen chimenea USD\,191\,000 frente a USD\,135\,000. En cambio, las viviendas remodeladas (año de remodelación posterior al de construcción) tienen una mediana menor (USD\,155\,000) que las no remodeladas (USD\,170\,000); esto no indica que remodelar reduzca el precio, sino que la variable binaria se confunde con la antigüedad de la vivienda, por lo que se analizará junto con el año de construcción (p.\,ej. años desde la remodelación).

Se detectaron dos valores atípicos (\texttt{Id} 524 y 1299): viviendas de más de 4\,000 pies$^2$ habitables y calidad 10 vendidas por menos de USD\,185\,000 (ambas con \texttt{SaleCondition} = \texttt{Partial}). Se eliminarán del entrenamiento.

\subsection{Configuración de aprendizaje y justificación}
El problema se aborda como \textbf{aprendizaje supervisado de regresión}: se dispone de ejemplos etiquetados (vivienda $\rightarrow$ precio real) y la salida es una variable continua. Se descartan configuraciones no supervisadas o por refuerzo porque el objetivo no es descubrir grupos ni tomar decisiones secuenciales, sino predecir un valor conocido en el histórico. Se compararán modelos lineales regularizados (Ridge, Lasso, Elastic Net) y modelos de ensamble (Random Forest, Gradient Boosting/XGBoost), evaluados con validación cruzada y con la métrica de la competencia, el RMSLE: $\mathrm{RMSLE}=\sqrt{\frac{1}{n}\sum_{i=1}^{n}\left(\log(1+\hat y_i)-\log(1+y_i)\right)^2}$, donde $y_i$ es el precio real y $\hat y_i$ el predicho; penaliza errores relativos, de modo que equivocarse en una casa barata pesa tanto como en una cara. También se reportarán $R^2$ y MAE.

%======================================================================
\section{Estado del arte}
Se revisaron cuatro trabajos que abordan el mismo problema (regresión del precio de venta) con la base Ames Housing. En todos el paradigma es \textbf{aprendizaje supervisado de regresión}.

\subsection{Sharma, Harsora y Ogunleye (2024) \cite{sharma2024}}
Artículo de revista (\textit{Analytics}). \textit{Técnicas:} regresión lineal, perceptrón multicapa, \textit{random forest}, SVR y XGBoost, con ajuste de hiperparámetros mediante \textit{GridSearchCV}, sobre la versión completa de Ames (2930 registros, 82 variables). \textit{Validación:} gráfica de valores reales vs. predichos e histograma de residuos del mejor modelo, además del puntaje de validación cruzada. \textit{Métricas:} $R^2$, $R^2$ ajustado ($1-(1-R^2)\frac{n-1}{n-k-1}$, con $k$ predictores), MSE, RMSE y MAE. \textit{Resultados:} XGBoost afinado fue el mejor ($R^2=0.930$, RMSE $=0.116$, MAE $=0.084$), frente a regresión lineal (RMSE $=0.130$, MAE $=0.075$) y \textit{random forest} (RMSE $=0.145$). Las variables más importantes fueron \texttt{OverallQual}, \texttt{GrLivArea}, \texttt{GarageCars} y \texttt{TotalBsmtSF}.

\subsection{Rath y Slater (2025) \cite{rath2025}}
Artículo de revista (\textit{SMU Data Science Review}). \textit{Técnicas:} regresión lineal múltiple, \textit{random forest} y XGBoost (100 árboles, tasa de aprendizaje 0.1, hiperparámetros fijos) sobre la base de Kaggle (79 variables explicativas). Imputan 15\,749 valores faltantes con media y moda, codifican las categóricas con \textit{label encoding} y estandarizan solo para la regresión lineal. \textit{Validación:} validación cruzada de 5 pliegues. \textit{Métricas:} RMSE, MAE y $R^2$ promedio, en dólares. \textit{Resultados:} XGBoost (RMSE 23\,274, MAE 14\,399, $R^2=0.913$) superó a \textit{random forest} (25\,439; 15\,712; 0.897) y a la regresión lineal (32\,362; 19\,501; 0.832). Las variables más influyentes fueron \texttt{OverallQual}, \texttt{GrLivArea}, \texttt{GarageCars} y las de sótano.

\subsection{Fourkiotis y Tsadiras (2023) \cite{fourkiotis2023}}
Artículo de congreso (AIAI 2023). \textit{Datos:} Ames de Kaggle (79 variables, 1458 viviendas), tras eliminar valores atípicos; aplican logaritmo al precio y a las variables asimétricas, imputan con ``None'', 0 o la moda, usan \textit{label encoding} y crean una variable de área total. \textit{Técnicas:} ElasticNet, Lasso, SVM, \textit{Kernel Ridge}, \textit{random forest}, \textit{Gradient Boosting}, XGBoost y LightGBM, más dos ensambles: \textit{Average} (promedio de seis modelos) y \textit{Voting} (SVR, XGB y LGBM con pesos 0.4, 0.4 y 0.2), con ajuste de hiperparámetros por \textit{GridSearch}. \textit{Validación:} validación cruzada de 5 pliegues y partición 80/20 (1166 para entrenamiento y 292 para prueba); reportan el desempeño en el conjunto de prueba. \textit{Métricas:} RMSLE (la de Kaggle) y $R^2$. \textit{Resultados:} \textit{Average} logró el mejor RMSLE (0.1071), seguido de \textit{Gradient Boosting} (0.1087) y \textit{Voting} (0.1093), que obtuvo el mejor $R^2$ (0.9193). Lasso y ElasticNet alcanzaron RMSLE de 0.1145 y 0.1146, mejor que \textit{random forest} (0.1325) y LightGBM (0.1348).

\subsection{Wang et al. (2026) \cite{wang2026}}
Artículo de revista (\textit{Science Journal of Applied Mathematics and Statistics}). Aborda el mismo problema (regresión del precio de vivienda) pero con otras bases de datos: Boston Housing (506 muestras, 13 variables) y California Housing (20\,640 observaciones, 8 variables). \textit{Técnicas:} OLS, Ridge, Lasso, Elastic Net y regresión polinómica de grado 2, con configuración por defecto ($\alpha=1.0$). \textit{Validación:} partición 80/20 entrenamiento/prueba, más análisis de residuos y gráficos Q-Q. \textit{Métricas:} MSE y $R^2$. \textit{Resultados:} la regresión polinómica fue la mejor (Boston: MSE 14.18, $R^2=0.81$; California: MSE $4.61\times10^{9}$, $R^2=0.66$), mientras OLS, Ridge, Lasso y Elastic Net rindieron de forma similar entre sí (Boston: $R^2\approx0.67$; California: $R^2=0.58$ en OLS, Ridge y Lasso).

\subsection{Síntesis}
Los trabajos coinciden en que los modelos más flexibles superan a la regresión lineal simple: los ensambles de árboles (XGBoost, \textit{random forest}) en \cite{sharma2024,rath2025}, los ensambles por promedio o votación en \cite{fourkiotis2023} y la expansión polinómica en \cite{wang2026}. Sobre Ames, la calidad general, el área habitable, el garaje y el sótano son los mayores determinantes del precio, lo cual concuerda con nuestro análisis exploratorio. En \cite{fourkiotis2023}, además, los modelos lineales regularizados (Lasso y ElasticNet) superaron a \textit{random forest} y LightGBM, por lo que PrediCasa comparará ambas familias. Las métricas no son directamente comparables: \cite{rath2025} reporta el error en dólares, \cite{sharma2024} valores cercanos a 0.1 que parecen calculados sobre el precio transformado y \cite{fourkiotis2023} el RMSLE; por eso PrediCasa reportará tanto RMSLE como error en dólares. Además, \cite{rath2025} incluye el identificador de parcela (\texttt{PID}) entre las variables importantes del \textit{random forest}, lo que sugiere fuga de información; PrediCasa descartará \texttt{Id}. Como limitación común, los datos provienen de una sola ciudad o región.

\begin{thebibliography}{9}
\bibitem{decock2011} D. De Cock, ``Ames, Iowa: Alternative to the Boston Housing Data as an End of Semester Regression Project,'' \textit{J. Stat. Educ.}, vol. 19, no. 3, 2011.
\bibitem{kaggle} Kaggle, ``House Prices: Advanced Regression Techniques.'' [En línea]. Disponible: \url{https://www.kaggle.com/c/house-prices-advanced-regression-techniques}
\bibitem{sharma2024} H. Sharma, H. Harsora y B. Ogunleye, ``An Optimal House Price Prediction Algorithm: XGBoost,'' \textit{Analytics}, vol. 3, no. 1, pp. 30--45, 2024, doi: 10.3390/analytics3010003.
\bibitem{rath2025} P. Rath y R. Slater, ``Predicting House Prices Using Machine Learning,'' \textit{SMU Data Science Review}, vol. 9, no. 3, art. 12, 2025.
\bibitem{fourkiotis2023} K. P. Fourkiotis y A. Tsadiras, ``Comparing Machine Learning Techniques for House Price Prediction,'' en \textit{Artificial Intelligence Applications and Innovations (AIAI 2023)}, IFIP AICT, vol. 676, Springer, 2023, pp. 292--303, doi: 10.1007/978-3-031-34107-6\_23.
\bibitem{wang2026} J. Wang \textit{et al.}, ``Unveiling the Impact of Nonlinear Modeling in Housing Price Prediction: An Empirical Comparative Study,'' \textit{Sci. J. Appl. Math. Stat.}, vol. 14, no. 1, pp. 6--15, 2026, doi: 10.11648/j.sjams.20261401.12.
\end{thebibliography}

\end{document}
